\documentclass[12pt]{article}
\usepackage[left=3cm, right=3cm, top=2.5cm]{geometry}  % set margins
\usepackage[document]{ragged2e}
\usepackage[utf8]{inputenc}
\usepackage[english]{babel}
\usepackage{graphicx}
\usepackage{hyperref}
\hypersetup{
  colorlinks=true,
  linkcolor=blue,
  filecolor=magenta,
  urlcolor=cyan,
  pdftitle={Hamish Morgan - Resume},
  pdfpagemode=FullScreen,
}
\setcounter{secnumdepth}{4}

%% \renewcommand{\familydefault}{\sfdefault}  % set font to this nice smooth one
%% \renewcommand{\familydefault}{\ttdefault}  % set font to this nice smooth one
%% \renewcommand{\baselinestretch}{1.5}

\begin{document}

\begin{center}
  \Large
  \vspace{0.8cm}
  \textbf{Hamish Morgan}\\
  \vspace{0.8cm}
  \large
  \href{https://github.com/hacmorgan}{GitHub @hacmorgan} - ham430@gmail.com - 0407245066

\end{center}


\begin{FlushLeft}

  \section*{Cover Letter}

  I have been working in various software roles in an AI \& robotics startup for seven years now. \\

  \vspace{0.3cm}
  
  In that time, I have seen the importance of good software engineering practices displayed over and over; from well-designed software empowering whole teams to work twice as efficiently, to unfavourable outcomes being reaped after questionable designs were sowed. I believe wholeheartedly in the value of good design and engineers who care about their tools, because I have seen how powerful they can be. \\

  \vspace{0.3cm}

  I personally love to build silly things that probably don't need to exist, and tinker with existing things, be they software or hardware (especially if I think they would work better with some modification). I find the best way to understand the world is to take it apart and look inside, and while this curiosity has occasionally cost me an electronic device or operating system installation, it has fostered an appreciation in me for the small details and ingenious solutions of the world.

\end{FlushLeft}


\end{document}
